\documentclass[11pt,a4paper]{article}
\usepackage[utf8]{inputenc}
\usepackage[T1]{fontenc}
\usepackage[margin=2.2cm]{geometry}
\usepackage{booktabs}
\usepackage{xcolor}
\usepackage{enumitem}
\usepackage{titlesec}
\usepackage{hyperref}
\usepackage{fancyhdr}

\definecolor{accent}{HTML}{1F6FEB}
\definecolor{ok}{HTML}{1A7F37}
\definecolor{warn}{HTML}{9A6700}
\hypersetup{colorlinks=true, urlcolor=accent, linkcolor=accent}

\titleformat{\section}{\large\bfseries\color{accent}}{}{0em}{}
\titlespacing{\section}{0pt}{10pt}{4pt}

\pagestyle{fancy}\fancyhf{}
\lhead{\small Dini Deploy}\rhead{\small Sprint Report}
\cfoot{\small\thepage}
\renewcommand{\headrulewidth}{0.4pt}

\newcommand{\done}{\textbf{\color{ok}Done}}
\newcommand{\prog}{\textbf{\color{warn}In Progress}}

\begin{document}

\begin{center}
{\LARGE\bfseries Weekly Sprint Report}\\[2pt]
{\large Deployment Model MultiDeepC ke Raspberry Pi 5 + Hailo-8L}\\[6pt]
{\small Sprint: Deployment Bring-up \quad|\quad Tanggal: \today \quad|\quad Owner: Riyadh}\\
{\small Repo: \url{https://github.com/riyadh-gif/dini-deploy} (private)}
\end{center}

\vspace{4pt}
\hrule
\vspace{8pt}

\section{Sprint Goal}
Mengemas model MultiDeepC (klasifikasi kondisi lahan per grid) menjadi container Docker
yang dapat dijalankan di perangkat edge Raspberry Pi 5 dengan akselerator Hailo-8L, lalu
memverifikasi seluruh jalur inferensi di hardware sebenarnya sebelum diserahkan ke tim hardware.

\section{Ringkasan Eksekutif}
Container Docker yang memakai NPU Hailo-8L sudah dibangun dan terbukti berjalan di Raspberry Pi 5.
Seluruh pipeline dari checkpoint PyTorch sampai inferensi di NPU sudah diverifikasi di hardware.
Latensi inferensi sekitar 0.3 ms per grid. Dua artefak preprocessing (scaler dan CNN) masih menjadi
dependensi terbuka untuk akurasi produksi.

\section{Pekerjaan Selesai Sprint Ini}
\begin{itemize}[leftmargin=1.4em,itemsep=2pt]
  \item \done\ Rekonstruksi arsitektur MultiDeepC dari checkpoint, strict load 115 tensor.
  \item \done\ Verifikasi clustering head DEC (Student-t, alpha 1.0) cocok output pelatihan (selisih 5e-7).
  \item \done\ Ekspor jalur encoder ke ONNX opset 11, seluruh operator didukung Hailo.
  \item \done\ Kompilasi HEF via Hailo Dataflow Compiler 3.33.1, target hailo8l.
  \item \done\ Build dan jalankan container Docker (backend CPU dan backend Hailo NPU) di Pi.
  \item \done\ Verifikasi di hardware: HEF load pada HailoRT 4.23.0, inferensi NPU end-to-end.
  \item \done\ Repo, README, dokumen serah terima, dan skrip PoC demo siap untuk tim hardware.
\end{itemize}

\section{Metrik Sprint}
\begin{center}
\begin{tabular}{@{}lll@{}}
\toprule
\textbf{Metrik} & \textbf{Nilai} & \textbf{Catatan} \\
\midrule
Latensi inferensi (NPU encode) & 0.27 ms median & per grid, sinkron \\
Latensi full path (NPU + DEC head) & 0.34 ms mean & termasuk clustering head CPU \\
Throughput streaming (NPU) & 15.480 inferensi/detik & mode pipelined \\
Latensi via container HTTP & 5 sampai 6 ms & termasuk FastAPI dan jaringan \\
Ukuran HEF & 762 KB & model terkompilasi \\
Image container Hailo & 403 MB & arm64, debian trixie \\
Operator tidak didukung Hailo & 0 & tidak perlu partisi graf \\
\bottomrule
\end{tabular}
\end{center}

\section{Demo}
\begin{enumerate}[leftmargin=1.4em,itemsep=1pt]
  \item \texttt{sudo docker ps} menampilkan container \texttt{mdc-hailo} berjalan di Pi.
  \item \texttt{GET /health} mengembalikan \texttt{"backend":"hailo"} (inferensi di NPU).
  \item \texttt{POST /predict} dengan data 1 grid mengembalikan label kondisi lahan.
  \item Skrip \texttt{poc\_demo.py} menampilkan tabel beberapa grid menuju label beserta latensi.
\end{enumerate}

\section{Blocker dan Risiko}
\begin{itemize}[leftmargin=1.4em,itemsep=2pt]
  \item \prog\ \textbf{Scaler pelatihan hilang.} Saat ini memakai placeholder identitas, sehingga nilai
        prediksi belum membedakan antar kelas. Perlu file scaler atau matrix input pelatihan.
  \item \prog\ \textbf{CNN feature extractor hilang.} Branch citra (64 fitur) diisi nol. Perlu model CNN
        beserta kode arsitekturnya untuk pipeline citra end-to-end.
  \item \textbf{Akurasi belum final.} Kompilasi memakai calibration sintetis. Setelah dua artefak di atas
        tersedia, dilakukan kompilasi ulang untuk akurasi produksi.
\end{itemize}

\section{Rencana Sprint Berikutnya}
\begin{itemize}[leftmargin=1.4em,itemsep=2pt]
  \item Memperoleh scaler dan CNN extractor dari pemilik model, lalu kompilasi ulang HEF.
  \item Validasi akurasi pasca kuantisasi terhadap output pelatihan.
  \item Menambahkan auto-start container di Pi (systemd atau restart policy) untuk operasi lapangan.
  \item Mengaktifkan jalur preprocessing on-device (perhitungan NDVI dan ekstraksi fitur citra).
\end{itemize}

\section{Tautan}
\begin{itemize}[leftmargin=1.4em,itemsep=1pt]
  \item Repositori (kode, artefak, dokumentasi): \url{https://github.com/riyadh-gif/dini-deploy}
  \item Panduan serah terima tim hardware: \texttt{docs/HANDOFF.md}
  \item Status preprocessing dan langkah lanjut: \texttt{docs/scaler\_status.md}
\end{itemize}

\end{document}
